\documentclass[11pt]{article}
\usepackage[margin=2.7cm]{geometry}
\usepackage{amsmath,amssymb,amsthm,mathtools}
\usepackage{enumitem}
\usepackage{booktabs}
\usepackage[colorlinks=true]{hyperref}
\usepackage[capitalize,noabbrev]{cleveref}

\newtheorem{theorem}{Theorem}[section]
\newtheorem{lemma}[theorem]{Lemma}
\newtheorem{corollary}[theorem]{Corollary}
\newtheorem{proposition}[theorem]{Proposition}
\theoremstyle{remark}
\newtheorem{remark}[theorem]{Remark}
\newtheorem{question}[theorem]{Open question}

\newcommand{\R}{\mathbb{R}}
\newcommand{\bu}{\boldsymbol{u}}
\newcommand{\bv}{\boldsymbol{v}}
\newcommand{\bw}{\boldsymbol{w}}
\newcommand{\ba}{\boldsymbol{a}}
\newcommand{\bff}{\boldsymbol{f}}
\newcommand{\bn}{\boldsymbol{n}}
\newcommand{\bz}{\boldsymbol{z}}
\newcommand{\ViscProj}{\Pi^{\mathrm{DS}}}
\newcommand{\norm}[1]{\lVert #1 \rVert}
\newcommand{\normK}[1]{\lVert #1 \rVert_{K}}
\newcommand{\tnorm}[1]{\lvert\!\lvert\!\lvert\, #1 \,\rvert\!\rvert\!\rvert}
\newcommand{\tns}{\tau_{1,\mathrm{NS}}}
\newcommand{\Eint}[2]{\mathrm{E}_{\mathrm{int},#1}(#2)}
\newcommand{\Dah}{\mathrm{Da}_{h,K}}
\newcommand{\Reh}{\mathrm{Re}_{h,K}}
\newcommand{\Pone}{\Pi_{1}}
\newcommand{\Ptwo}{\Pi_{2}}
\newcommand{\Poneo}{\Pi_{1}^{\perp}}
\newcommand{\Ptwoo}{\Pi_{2}^{\perp}}
\newcommand{\Xhz}{\mathcal{X}_{h0}}
\newcommand{\st}{\widetilde{\sigma}_\alpha}
\newcommand{\Bs}{B_{\mathrm{S}}}
\newcommand{\Bo}{B_{\mathrm{osgs}}}
\newcommand{\EF}{\mathcal{E}}

\title{Can the ASGS and OSGS convergence results share one proof?\\[1ex]
\large An exploration: a provable dichotomy, a unified penalized theorem,\\
and what should stay method-specific}
\author{}
\date{\today}

\begin{document}
\maketitle

\begin{abstract}
We examine whether the inf--sup proof strategy developed for the OSGS variant
in~\cite{OSGSnote} can be transported to the ASGS variant of~\cite{companion}, so
that both a priori convergence results follow from essentially one argument. The
answer has a sharp boundary. \emph{Negative part:} full-strength reactive control
$\norm{\sigma^{1/2}\bu_h}$ is provably unattainable for the ASGS form---the
stabilization renormalizes the reaction to $\st$ against \emph{every} test function,
and we exhibit a family of discrete velocities along which the ASGS inf--sup constant
in the strong (OSGS) norm degenerates like $(1+\Dah)^{-1/2}$. There is therefore no
common norm, and the $\sigma$-versus-$\st$ dichotomy identified in~\cite{OSGSnote} is
a property of the methods, not of the proofs. \emph{Positive part:} everything short
of the norm does unify. We prove a jump-condition-free convergence theorem for the
ASGS method in its own triple norm, with the same error function and the same
$(1+\Dah)^{1/2}$ factor as the OSGS theorem, by the same skeleton: shared parameter
inequalities, a shared interpolation core, exact consistency replacing the
best-approximation consistency estimate, two zero-cost telescopings inherited from
the companion's own continuity argument, and a short inventory of stabilization
products of which exactly three families carry the Damk\"ohler factor. The two Damk\"ohler-sensitive results are then
quantitatively symmetric: ASGS concedes $\Dah^{1/2}$ in the \emph{norm}
($\st/\sigma \sim \Dah^{-1}$), OSGS concedes it in the \emph{constant}. The
companion's Theorem~5.3 remains the sharper ASGS statement---its Damk\"ohler-freeness
is purchased by the facewise $[\tau_1]$-estimates spent on the face terms of its
integration-by-parts telescopings---and we recommend keeping it. Finally we examine
the jump condition itself (\cref{sec:jump}): it is a facewise comparability of
$\varphi_1 = \alpha_K\tns^{-1}$, \emph{implied} by the standing hypotheses with
$h$-uniform constants (the same content, established the same way, as the
patch-equivalence lemma on which the OSGS analysis rests), and it has \emph{no point
of application} in the OSGS analysis, whose orthogonal projection annihilates the
very pairings whose elementwise integration by parts the condition prices. A
best-approximation floor (\cref{prop:floor}) shows that no hypothesis and no method
can remove the $(1+\Dah)^{1/2}$ factor from a strong-norm estimate. The hypothesis
ledger is thus rebalanced: neither the jump condition nor the enlarged inverse
constant is a genuine surcharge of the ASGS analysis, and the only assumption in
either proof not implied by the standing data and mesh hypotheses is the
projection-stability condition (A7), on the OSGS side.
\end{abstract}

\section{The question and the verdict}
\label{sec:verdict}

The companion manuscript~\cite{companion} proves stability and convergence for the
ASGS variant by coercivity plus a continuity lemma (with an interior-face jump
condition on $\tau_1$ when $\sigma>0$); the note~\cite{OSGSnote} proves the analogous
results for the OSGS variant by the inf--sup strategy of~\cite{Codina2008}
(special test function, projection-stability hypothesis, best-approximation
consistency). The natural question: can the second strategy absorb the first, so that
both theorems are two instantiations of one proof?

The exploration below yields a three-part verdict.

\begin{enumerate}[label=(\Roman*), leftmargin=2.5em, itemsep=0.4ex]
\item \emph{The proof skeleton unifies; the theorems do not.} Roughly $80\%$ of the
material is genuinely shared: hypotheses, parameter inequalities, the interpolation
core, the error-function format, and the overall
stability--consistency--interpolation architecture. What cannot be shared is the
norm. \Cref{prop:negative} shows that the strong norm of~\cite{OSGSnote} (with full
$\norm{\sigma^{1/2}\bu_h}$) is unattainable for the ASGS bilinear form: the inf--sup
constant in that norm degenerates like $(1+\Dah)^{-1/2}$. The reactive
renormalization $\sigma \mapsto \sigma(1-\sigma\tau_1) = \st$ is an identity of the
bilinear form, active against every test function, so no test-function ingenuity can
undo it. Symmetrically, the OSGS form cannot produce the $\st$-renormalization
(its projection deletes the reactive residual), so it keeps the strong norm and pays
$(1+\Dah)^{1/2}$ in the constant~\cite{OSGSnote}. The dichotomy is a theorem about
the methods.
\item \emph{A unified, jump-free, penalized theorem holds.}
\Cref{th:unifiedASGS} proves for the ASGS method, in its own triple norm and
\emph{without} the jump condition of~\cite{companion}, the estimate
\[
  \tnorm{U - U_h}_{\mathrm{S}}
  \;\le\;
  C \sum_K \frac{\alpha_K}{h_K}\,\bigl(1+\Dah\bigr)^{1/2}
  \Bigl( \tau_{2,K}^{1/2}\Eint{K}{\bu} + \tau_{1,K}^{1/2}\Eint{K}{p} \Bigr),
\]
which is the OSGS theorem of~\cite{OSGSnote} verbatim, norm apart. Its proof reuses
the interpolation core of~\cite{OSGSnote} (three of five Galerkin estimates
unchanged, the remaining two entering zero-cost telescopings with their
stabilization partners), replaces the best-approximation consistency lemma by exact
consistency, and adds a short inventory of stabilization products, of which exactly
three families carry the Damk\"ohler factor. ``Jump-free'' is here a statement about
the proof, not a hypothesis saving: \cref{sec:jump} shows the jump condition costs
nothing to assume.
\item \emph{The sharp ASGS theorem should be kept.} The Damk\"ohler-free Theorem~5.3
of~\cite{companion} is strictly stronger than \cref{th:unifiedASGS} in its regime,
and its extra machinery (the $\tau_1$ jump condition, used in an elementwise
integration by parts) is spent on the pressure and convective \emph{face terms} of
its two integration-by-parts telescopings, and on nothing else. \Cref{sec:jump}
shows that the underlying facewise comparability is implied by the standing
hypotheses, so the sharp theorem's hypothesis set is effectively the standing one;
\cref{rem:sharper} sketches an alternative, regime-splitting route to the same face
terms. Our recommendation (\cref{sec:recommendation}) remains to keep both existing
proofs and add a remark on the shared architecture, the dichotomy, and the corrected
hypothesis ledger.
\item \emph{The jump condition is free for ASGS and inert for OSGS.}
\Cref{sec:jump} examines the hypothesis itself. It is a facewise comparability
requirement on $\varphi_1 = \alpha_K\tns^{-1}$ ($\sigma$ cancels in the jump of
$\tau_1^{-1} = \varphi_1 + \sigma$), and \cref{lem:jumpauto} shows it is implied,
with $h$-uniform constants, by the standing mesh, porosity and advection
hypotheses---by the same mechanism, with the same $\nu$-floor, as the
patch-equivalence and smoothing lemmas that the OSGS analysis itself consumes.
Conversely, granting it to the OSGS analysis improves nothing: its only role is to
price the interior-face terms generated by elementwise integration by parts of the
$\sigma$-weighted stabilization pairings, and the OSGS projection annihilates those
pairings, so no such face term ever arises (\cref{rem:inert}); moreover the dominant
OSGS penalty is forced by the norm at the level of best approximation
(\cref{prop:floor}), hence unremovable by any hypothesis and any method. The net
hypothesis ledger favors, if anything, the ASGS analysis
(\cref{rem:ledger}).
\end{enumerate}

An honest note on the exploration path: our first attempt was to \emph{restore} full
$\sigma$-control for ASGS by enriching the special test function with a component
$\tau_1\circledast(\sigma\bu_h)$, whose Galerkin pairing produces exactly the missing
$\sum_K \sigma^2\tau_{1,K}\norm{\bu_h}_K^2$ (note the exact identity
$\sigma^2\tau_1 = \sigma - \st$). This fails, instructively: the stabilization's
reactive--adjoint coupling returns $-\sigma^2\tau_1(\bw,\bu_h)$ against the same
component, cancelling the gain to leading order---the renormalization
$\sigma(1-\sigma\tau_1)=\st$ applies to the pair $(\bu_h,\bw)$ for \emph{any} $\bw$.
That observation, made quantitative, is \cref{prop:negative}.

\section{Setting}
\label{sec:setting}

We work in the setting, notation and hypotheses (A1)--(A8) of~\cite{OSGSnote}, which
restates~\cite{companion}; we recall only what is used here. The linearized problem,
the Galerkin form $B$, the composite $X(V) = \ba\cdot\nabla\bv + \nabla q$, the
parameters $\tau_{1,K}, \tau_{2,K}, \tns$, the elementwise numbers
$\Dah = \sigma h_K^2/(\alpha_K\nu)$ and $\Reh$, the weighted products
$(\cdot,\cdot)_{\tau_i}$, the projections $\Pone, \Ptwo$ onto the unconstrained
spaces with fluctuations $\Poneo, \Ptwoo$, and the interpolation quantities
$\Eint{K}{\cdot}$ are all as there. The parameter inequalities (P1)--(P6) and
(K1)--(K3) of~\cite[Lemma~4.1]{OSGSnote} are used freely. We also use
$\delta_\alpha \coloneqq \max_K \alpha_{\infty,K}/\alpha_{0,K}$, bounded under the
porosity-resolution condition, and $M_2 \coloneqq 1 + C_\nabla/C_{\mathrm{inv}}$.

The two stabilized forms are, with $G(\bu) \coloneqq
2\nabla\cdot(\alpha\nu\ViscProj\nabla\bu) - \sigma\bu$,
\begin{align}
  \Bs(\ba; U, V)
  &\coloneqq B(\ba; U, V)
  + \sum_K \tau_{1,K}\bigl( \alpha X(V) + G(\bv),\ \alpha X(U) - G(\bu) \bigr)_K
  \notag\\
  &\qquad\qquad
  + \sum_K \tau_{2,K}\bigl( \nabla\cdot(\alpha\bv) - \varepsilon q,\
  \varepsilon p + \nabla\cdot(\alpha\bu) \bigr)_K,
  \label{eq:BSdef}\\
  \Bo(\ba; U, V)
  &\coloneqq B(\ba; U, V)
  + \bigl( \Poneo[\alpha X(U)],\, \alpha X(V) \bigr)_{\tau_1}
  + \bigl( \Ptwoo[\nabla\cdot(\alpha\bu)],\, \nabla\cdot(\alpha\bv) \bigr)_{\tau_2}.
  \label{eq:Bodef}
\end{align}
The diagonal of \cref{eq:BSdef} reproduces the stability identity
of~\cite{companion} (the crossed viscous--reactive products telescope into
$\norm{\tau_1^{1/2}\alpha X(U)}_h^2 - \norm{\tau_1^{1/2}G(\bu)}_h^2$ and
$\norm{\tau_2^{1/2}\nabla\cdot(\alpha\bu)}_h^2 - \varepsilon^2
\norm{\tau_2^{1/2}p}_h^2$), so \cref{eq:BSdef} is the ASGS form of the companion
written out; \cref{eq:Bodef} is the analyzed OSGS form of~\cite{OSGSnote}. The two
triple norms are
\begin{align}
  \tnorm{V}^2
  &\coloneqq
  2\nu\norm{\alpha^{1/2}\ViscProj\nabla\bv}^2
  + \norm{\sigma^{1/2}\bv}^2
  + \varepsilon\norm{q}^2
  + \norm{\alpha X(V)}_{\tau_1}^2
  + \norm{\nabla\cdot(\alpha\bv)}_{\tau_2}^2,
  \label{eq:StrongNorm}\\
  \tnorm{V}_{\mathrm{S}}^2
  &\coloneqq
  \nu\norm{\alpha^{1/2}\ViscProj\nabla\bv}^2
  + \norm{\st^{1/2}\bv}_h^2
  + \varepsilon\norm{q}^2
  + \norm{\alpha X(V)}_{\tau_1}^2
  + \norm{\nabla\cdot(\alpha\bv)}_{\tau_2}^2,
  \label{eq:ASGSNorm}
\end{align}
where $\st \coloneqq \tns^{-1}\sigma/(\tns^{-1}+\sigma/\alpha_K)$, and
$\tnorm{\cdot}_{\mathrm{S}} \le \tnorm{\cdot}$ termwise. Three exact identities
organize everything (elementwise; direct algebra from
$\tau_1^{-1} = \alpha_K\tns^{-1} + \sigma$):
\begin{equation}
  \label{eq:ThreeIdentities}
  \sigma(1 - \sigma\tau_1) = \st,
  \qquad
  \sigma^2\tau_1 = \sigma - \st,
  \qquad
  \frac{\sigma\tau_1}{1-\sigma\tau_1} = \frac{\sigma\tns}{\alpha_K} \le \frac{\Dah}{c_1}.
\end{equation}
The first says the ASGS stabilization renormalizes the reaction; the second
identifies the renormalization deficit; the third is the exchange rate for
converting undamped reactive quantities into $\st$-quantities, and is the source of
every Damk\"ohler factor below. Note also $\st \le \alpha_K\tns^{-1} =
c_1\alpha_K^2\tau_{2,K}/h_K^2$, i.e.,
\begin{equation}
  \label{eq:SigmaTildeScale}
  \st^{1/2} \;\le\; \sqrt{c_1}\,\frac{\alpha_K}{h_K}\,\tau_{2,K}^{1/2},
\end{equation}
so $\st$-weighted quantities are always commensurate with the
divergence-stabilization scale: this is why the renormalized reaction never costs a
Damk\"ohler factor.

\section{The dichotomy is a theorem: no common norm}
\label{sec:dichotomy}

The OSGS analysis of~\cite{OSGSnote} controls the full reactive term
$\norm{\sigma^{1/2}\bu_h}$; the ASGS analysis of~\cite{companion} controls only
$\norm{\st^{1/2}\bu_h}_h$. If the second were an artifact of the coercivity-based
proof, an inf--sup argument with a cleverer test function might upgrade it, and both
methods could then be analyzed in the single strong norm \cref{eq:StrongNorm}. The
following proposition rules this out. To keep the counterexample free of inessential
technicalities we place it in the Darcy-dominated corner of the parameter space,
which is precisely the regime where the two norms differ.

\begin{proposition}[Sharpness of the reactive renormalization]
\label{prop:negative}
Let $\ba \equiv \boldsymbol{0}$, $\alpha \equiv 1$, $\varepsilon = 0$, let
$\sigma, \nu > 0$ be constant, and let the mesh be such that $\tau_{1,K}$ and
$\tau_{2,K}$ take a single value on all elements (e.g., a uniform mesh), with the
spaces of~\cite{OSGSnote} and $\dim V_{h0} > \dim Q_h - 1$ (true for every standard
pair once the mesh is fine enough). Write $\mathrm{Da}_h = \sigma h^2/\nu$
(so $\sigma\tau_1 = \mathrm{Da}_h/(c_1 + \mathrm{Da}_h)$ and
$\st/\sigma = c_1/(c_1+\mathrm{Da}_h)$). Then there exists
$\bu_h \in V_{h0}\setminus\{\boldsymbol 0\}$ such that, with
$U_h \coloneqq [\bu_h; 0]$,
\begin{equation}
  \label{eq:NegativeBound}
  \sup_{V_h \in \Xhz}
  \frac{\Bs(\ba; U_h, V_h)}{\tnorm{V_h}}
  \;\le\;
  C\,\sigma^{1/2}\bigl(1 + \mathrm{Da}_h\bigr)^{-1/2}\,\norm{\bu_h}
  \;\le\;
  C\,\bigl(1 + \mathrm{Da}_h\bigr)^{-1/2}\,\tnorm{U_h},
\end{equation}
with $C$ depending only on $c_1$, $c_2$ and $C_{\mathrm{inv}}$. Consequently the
inf--sup constant of $\Bs$ with respect to the strong norm \cref{eq:StrongNorm}
degenerates at least like $(1+\mathrm{Da}_h)^{-1/2}$ in the reaction-dominated
regime, and no stability estimate of the form
$\sup_V \Bs(U_h,V)/\tnorm{V} \ge \beta\tnorm{U_h}$ can hold with $\beta$ independent
of $\mathrm{Da}_h$.
\end{proposition}

\begin{proof}
\emph{Choice of $\bu_h$.} Since $\dim V_{h0} > \dim \nabla Q_h \le \dim Q_h - 1$,
there is $\bu_h \in V_{h0}$, $\norm{\bu_h} = 1$, with
$(\nabla q_h, \bu_h) = 0$ for all $q_h \in Q_h$; equivalently, since
$\bu_h \in H^1_0(\Omega)^d$, $(q_h, \nabla\cdot\bu_h) = 0$ for all $q_h \in Q_h$.

\emph{Expansion.} With $\ba = \boldsymbol 0$, $\alpha \equiv 1$, $p_h = 0$:
$X(U_h) = \boldsymbol 0$, $G(\bu_h) = 2\nu\nabla\cdot(\ViscProj\nabla\bu_h)
- \sigma\bu_h$, and for $V_h = [\bv_h; q_h]$,
\begin{align}
  \Bs(\ba; U_h, V_h)
  ={}& 2\nu(\ViscProj\nabla\bv_h, \ViscProj\nabla\bu_h)
  + \sigma(\bv_h, \bu_h)
  + (q_h, \nabla\cdot\bu_h)
  \notag\\
  &- \sum_K \tau_1\bigl( \nabla q_h + G(\bv_h),\ G(\bu_h) \bigr)_K
  + \sum_K \tau_2\bigl( \nabla\cdot\bv_h,\ \nabla\cdot\bu_h \bigr)_K .
  \label{eq:NegExpansion}
\end{align}
The third term vanishes by the choice of $\bu_h$. In the fourth, expand
$G(\bu_h)$ and $G(\bv_h)$ and collect the six products; two of them are special.
First, the reactive--pressure product: since $\tau_1$ is globally constant,
$+\sigma\tau_1(\nabla q_h, \bu_h) = \sigma\tau_1 \cdot 0 = 0$. Second, the
reactive--reactive product $-\sigma^2\tau_1(\bv_h,\bu_h)$ telescopes with the
Galerkin reaction by the first identity in \cref{eq:ThreeIdentities}:
\[
  \sigma(\bv_h,\bu_h) - \sigma^2\tau_1(\bv_h,\bu_h)
  = \st\,(\bv_h,\bu_h)
  \;\le\; \st^{1/2}\norm{\bv_h}\cdot \st^{1/2}\norm{\bu_h}
  \;\le\; \Bigl(\frac{c_1}{c_1+\mathrm{Da}_h}\Bigr)^{1/2}\sigma^{1/2}
  \norm{\bu_h}\,\tnorm{V_h},
\]
using $\st^{1/2}\norm{\bv_h} \le \sigma^{1/2}\norm{\bv_h} \le \tnorm{V_h}$. Every
remaining term carries at least one viscous factor on the $\bu_h$ side and is
Damk\"ohler-damped; we list the bounds, writing
$\nu^{1/2}/h = (\sigma/\mathrm{Da}_h)^{1/2}$ and using the inverse estimate
$\norm{\nabla\cdot(\ViscProj\nabla\bu_h)}_K \le (C_{\mathrm{inv}}/h)
\norm{\ViscProj\nabla\bu_h}_K \le (C_{\mathrm{inv}}/h)^2\norm{\bu_h}_K$ throughout:
{\small
\begin{align*}
  2\nu(\ViscProj\nabla\bv_h, \ViscProj\nabla\bu_h)
  &\le \sqrt 2\,\tnorm{V_h}\cdot \sqrt 2\,C_{\mathrm{inv}}
  \nu^{1/2}h^{-1}\norm{\bu_h}
  = 2C_{\mathrm{inv}}\,\mathrm{Da}_h^{-1/2}\,\sigma^{1/2}
  \norm{\bu_h}\tnorm{V_h},\displaybreak[0]\\
  2\nu\tau_1\bigl|(\nabla q_h, \nabla\cdot(\ViscProj\nabla\bu_h))\bigr|
  &\le \norm{\nabla q_h}_{\tau_1}\cdot
  2\nu\tau_1^{1/2}C_{\mathrm{inv}}^2 h^{-2}\norm{\bu_h}\notag\\
  &\le 2C_{\mathrm{inv}}^2\,\mathrm{Da}_h^{-1}\,\sigma^{1/2}\norm{\bu_h}\tnorm{V_h},
  \displaybreak[0]\\
  2\nu\sigma\tau_1\bigl|(\nabla\cdot(\ViscProj\nabla\bv_h), \bu_h)\bigr|
  &\le \sigma\tau_1\cdot 2\nu^{1/2}C_{\mathrm{inv}}h^{-1}\cdot
  \nu^{1/2}\norm{\ViscProj\nabla\bv_h}\,\norm{\bu_h}\notag\\
  &\le \frac{2C_{\mathrm{inv}}\mathrm{Da}_h^{1/2}}{c_1+\mathrm{Da}_h}\,
  \sigma^{1/2}\norm{\bu_h}\tnorm{V_h},\\
  4\nu^2\tau_1\bigl|(\nabla\cdot(\ViscProj\nabla\bv_h),
  \nabla\cdot(\ViscProj\nabla\bu_h))\bigr|
  &\le 4 C_{\mathrm{inv}}^3\,\mathrm{Da}_h^{-1}\,
  \sigma^{1/2}\norm{\bu_h}\tnorm{V_h},\\
  2\nu\sigma\tau_1\bigl|(\bv_h, \nabla\cdot(\ViscProj\nabla\bu_h))\bigr|
  &\le \sigma\tau_1\cdot 2\nu C_{\mathrm{inv}}^2h^{-2}
  \,\sigma^{-1/2}\cdot\sigma^{1/2}\norm{\bv_h}\norm{\bu_h}\notag\\
  &\le 2C_{\mathrm{inv}}^2\,\mathrm{Da}_h^{-1}\,
  \sigma^{1/2}\norm{\bu_h}\tnorm{V_h},\\
  \tau_2\bigl|(\nabla\cdot\bv_h, \nabla\cdot\bu_h)\bigr|
  &\le \norm{\nabla\cdot\bv_h}_{\tau_2}\cdot
  \tau_2^{1/2}\sqrt d\,C_{\mathrm{inv}}h^{-1}\norm{\bu_h}\notag\\
  &\le C\,\mathrm{Da}_h^{-1/2}\sigma^{1/2}\norm{\bu_h}\tnorm{V_h},
\end{align*}}%
where in the second line $\nu\tau_1^{1/2}h^{-2} \le \nu^{1/2}\sigma^{-1/2}
\nu^{1/2}h^{-2} = \sigma^{1/2}\mathrm{Da}_h^{-1}$ (via $\tau_1 \le 1/\sigma$), in
the third $\sigma\tau_1\,\nu^{1/2}h^{-1} =
\bigl[\mathrm{Da}_h/(c_1+\mathrm{Da}_h)\bigr]\mathrm{Da}_h^{-1/2}\sigma^{1/2}$, and
in the last $\tau_2 = (\nu)(1 + (c_2/c_1)\Reh)$ with $\Reh = 0$ here since
$\ba=\boldsymbol 0$. Summing all contributions,
\[
  \bigl| \Bs(\ba; U_h, V_h) \bigr|
  \le C\,\sigma^{1/2}\Bigl[
  \Bigl(\tfrac{c_1}{c_1+\mathrm{Da}_h}\Bigr)^{1/2}
  + \mathrm{Da}_h^{-1/2}
  + \tfrac{\mathrm{Da}_h^{1/2}}{c_1+\mathrm{Da}_h}
  + \mathrm{Da}_h^{-1}
  \Bigr]\norm{\bu_h}\,\tnorm{V_h},
\]
and every bracketed term is $\le C'(1+\mathrm{Da}_h)^{-1/2}$ for
$\mathrm{Da}_h \ge 1$ (for $\mathrm{Da}_h \le 1$ the claim
\cref{eq:NegativeBound} is trivial with $C$ adjusted). Finally
$\tnorm{U_h} \ge \norm{\sigma^{1/2}\bu_h} = \sigma^{1/2}\norm{\bu_h}$, which gives
the second inequality in \cref{eq:NegativeBound}.
\end{proof}

\begin{remark}[Reading of \cref{prop:negative}, and relation to the companion's
remark]
\label{rem:reading}
The companion manuscript remarks, after its stability lemma, that the weaker
apparent control of $\bu_h$ through $\st$ for large reaction ``turns out not to be
the case'', pointing to the robustness analysis and the numerics.
\Cref{prop:negative} does not contradict that statement; it sharpens the sense in
which it is true. What is genuinely weaker is the \emph{norm control}: no argument
can extract $\norm{\sigma^{1/2}\bu_h}$ from $\Bs$ with a Damk\"ohler-uniform
constant, because $\Bs$ simply does not contain that information along
gradient-orthogonal velocity directions in the Darcy limit. What is nevertheless
unaffected is the \emph{convergence rate}: the error equation only requires control
along the error direction, and the companion's Theorem~5.3 delivers optimal,
Damk\"ohler-free rates in the $\st$-norm. Both facts are correct simultaneously; the
proposition merely certifies that the $\st$-damping in the stability lemma is sharp,
not an estimate that a better proof might remove.
\end{remark}

\begin{remark}[Quantitative symmetry of the dichotomy]
\label{rem:symmetry}
In the Darcy limit $\st/\sigma = c_1/(c_1+\mathrm{Da}_h)$, so the ASGS norm concedes
a factor $\sim\Dah^{1/2}$ in its reactive component relative to the strong norm.
The OSGS theorem of~\cite{OSGSnote} keeps the strong norm and concedes the factor
$(1+\Dah)^{1/2}$ in the error constant. The two methods thus pay the same
Damk\"ohler toll on opposite sides of the inequality:
\[
  \underbrace{\tnorm{E_h}_{\mathrm{S}} \le C\,\EF_0(h)}_{\text{ASGS, sharp
  (Thm.~5.3 of~\cite{companion})}}
  \qquad\text{versus}\qquad
  \underbrace{\tnorm{E_h} \le C\,(1+\mathrm{Da}_h)^{1/2}\,\EF_0(h)}_{\text{OSGS
  (\cite{OSGSnote})}},
\]
with $\EF_0(h) \coloneqq \sum_K (\alpha_K/h_K)(\tau_{2,K}^{1/2}\Eint{K}{\bu} +
\tau_{1,K}^{1/2}\Eint{K}{p})$ the common error function, and, by
\cref{prop:negative}, neither statement can be moved to the other side's norm
without loss. This is the precise sense in which the methods---not the proofs---
differ.
\end{remark}

\begin{remark}[The deficit table and the test-function recipe]
\label{rem:deficits}
The unifying principle behind both stability proofs is: \emph{the special test
function must supply exactly the residual components that the method's own
stabilization leaves uncontrolled at the diagonal}. The deficits are:
\begin{center}
\begin{tabular}{lcc}
\toprule
diagonal control of & ASGS \cref{eq:BSdef} & OSGS \cref{eq:Bodef} \\
\midrule
first-order residual $\alpha X(U_h)$ & full & fluctuation only
($\Poneo$-part) \\
divergence $\nabla\cdot(\alpha\bu_h)$ & full & fluctuation only
($\Ptwoo$-part) \\
reaction & damped: $\sigma(1-\sigma\tau_1) = \st$ & full $\sigma$ \\
compressibility & damped: $\varepsilon(1-\varepsilon\tau_2)$ & full $\varepsilon$ \\
\bottomrule
\end{tabular}
\end{center}
For OSGS the missing pieces are the resolvable parts $\Pone[\alpha X(U_h)]$,
$\Ptwo[\nabla\cdot(\alpha\bu_h)]$, and the test components
$\tau_1\circledast\Pi_{1,0}(\cdot)$, $\tau_2\circledast\Ptwo(\cdot)$
of~\cite{OSGSnote} supply them (at the price of hypothesis (A7)). For ASGS the
missing pieces are the renormalization deficits $\sigma^2\tau_1\norm{\bu_h}^2_K$ and
$\varepsilon^2\tau_2\norm{p_h}^2_K$ (second identity in \cref{eq:ThreeIdentities}),
and \cref{prop:negative} shows that the reactive one \emph{cannot} be supplied: the
candidate component $\tau_1\circledast(\sigma\bu_h)$ generates its own gain and its
own cancellation in equal measure. The compressibility deficit is benign as long as
$\varepsilon\tau_2 \le C_2 < 1$, which the companion assumes. The recipe therefore
unifies the two stability proofs conceptually, with the ASGS instantiation
degenerating to pure coercivity (empty test enrichment) in its natural norm
\cref{eq:ASGSNorm}---which is the companion's existing lemma.
\end{remark}

\section{The unified penalized theorem: jump-free ASGS convergence}
\label{sec:unified}

We now prove the positive half of the verdict: run the convergence half of the
OSGS analysis of~\cite{OSGSnote} on the ASGS form, each method in its own norm.
The result is the OSGS theorem with $\tnorm{\cdot}$ replaced by
$\tnorm{\cdot}_{\mathrm{S}}$, proved without the interior-face jump condition
of~\cite{companion}. Throughout this section we assume (A1)--(A6), (A8)
of~\cite{OSGSnote} together with the hypotheses of the ASGS stability lemma
of~\cite{companion} (in particular $c_1 > 2\xi C_{\mathrm{inv}}^2$ with $\xi > 2$,
where $C_{\mathrm{inv}}$ is the enlarged constant of the companion's weighted
inverse estimate, controlled by $\delta_\alpha^{1/2}M_2$ times the plain constant);
hypothesis (A7) is \emph{not} needed (no projection appears), and neither is the
jump condition---though, as \cref{sec:jump} shows, avoiding the latter saves
nothing, since it is implied by the standing hypotheses; the value of the theorem
below is the shared architecture at a quantified price, not hypothesis economy.

The single method-specific algebraic ingredient is the telescoping identity, which
we isolate:

\begin{lemma}[Reactive telescoping]
\label{lem:telescope}
For any $\bw, \bz \in L^2(\Omega)^d$,
\begin{equation}
  \label{eq:Telescope}
  \sigma(\bz, \bw) - \sum_K \sigma^2\tau_{1,K}(\bz,\bw)_K
  = \sum_K \st\big|_K\, (\bz, \bw)_K ,
\end{equation}
and, elementwise, $\st^{1/2} \le \sqrt{c_1}\,(\alpha_K/h_K)\,\tau_{2,K}^{1/2}$.
\end{lemma}

\begin{proof}
The first claim is the first identity of \cref{eq:ThreeIdentities} applied
elementwise ($\tau_1$ is elementwise constant); the second is
\cref{eq:SigmaTildeScale}.
\end{proof}

\begin{lemma}[Jump-free interpolation continuity for the ASGS form]
\label{lem:asgsinterp}
Let $\hat E = U - \hat U_h$ be the interpolation error of~\cite[\S 6.2]{OSGSnote}.
Under the stated hypotheses, for all $V_h \in \Xhz$,
\begin{equation}
  \label{eq:ASGSInterpBound}
  \bigl| \Bs(\ba; \hat E, V_h) \bigr|
  \;\le\; C\,\EF(h)\,\tnorm{V_h}_{\mathrm{S}},
  \qquad
  \EF(h) = \sum_K \frac{\alpha_K}{h_K}(1+\Dah)^{1/2}
  \bigl( \tau_{2,K}^{1/2}\Eint{K}{\bu} + \tau_{1,K}^{1/2}\Eint{K}{p} \bigr).
\end{equation}
\end{lemma}

\begin{proof}
Split $\Bs(\hat E, V_h)$ into the Galerkin part and the stabilization part.

\smallskip\noindent\emph{Galerkin part, with two telescopings.}
Estimates (I1) (viscous), (I4) (compressibility) and (I5) (pressure gradient, by
global integration by parts) of~\cite[Lemma~6.2]{OSGSnote} apply verbatim: each uses
only the $\nu$-, $\varepsilon$- and
$\norm{\nabla\cdot(\alpha\bv_h)}_{\tau_2}$-components of the test norm, all present
in $\tnorm{\cdot}_{\mathrm{S}}$, and each is bounded by
$C\,\EF(h)\tnorm{V_h}_{\mathrm{S}}$, the Damk\"ohler factor now arising only in the
(I4)/(I5) conversions through (P5)-type inequalities of~\cite[Lemma~4.1]{OSGSnote}. The two remaining
Galerkin pairings are treated \emph{jointly with their stabilization partners},
exploiting the elementwise identity $1-\sigma\tau_1 = \varphi_1\tau_1$,
$\varphi_1 \coloneqq \alpha_K\tns^{-1}$---the mechanism of the companion's own
continuity proof. First, the Galerkin reaction telescopes with the
reactive--reactive stabilization product via \cref{lem:telescope}:
\[
  \sigma(\hat{\boldsymbol e}_u, \bv_h)
  - \sum_K \sigma^2\tau_{1,K}(\bv_h, \hat{\boldsymbol e}_u)_K
  = \sum_K \st(\hat{\boldsymbol e}_u, \bv_h)_K
  \le \Bigl( \sum_K c_1 \Bigl(\tfrac{\alpha_K}{h_K}\Bigr)^2
  \tau_{2,K}\,\Eint{K}{\bu}^2 \Bigr)^{1/2} \tnorm{V_h}_{\mathrm{S}},
\]
using $\st^{1/2}\normK{\hat{\boldsymbol e}_u} \le
\sqrt{c_1}(\alpha_K/h_K)\tau_{2,K}^{1/2}\Eint{K}{\bu}$ and
$\norm{\st^{1/2}\bv_h}_h \le \tnorm{V_h}_{\mathrm{S}}$; no Damk\"ohler factor
arises. Second, the convective and mass--divergence Galerkin terms, combined by
global integration by parts into $-(\hat{\boldsymbol e}_u, \alpha X(V_h))$ exactly
as in (I2) of~\cite{OSGSnote}, telescope with the stabilization product
$\sum_K \sigma\tau_{1,K}(\alpha X(V_h), \hat{\boldsymbol e}_u)_K$, which we
accordingly remove from the inventory below:
\[
  -\bigl(\hat{\boldsymbol e}_u, \alpha X(V_h)\bigr)
  + \sum_K \sigma\tau_{1,K}\bigl(\alpha X(V_h), \hat{\boldsymbol e}_u\bigr)_K
  = -\sum_K \varphi_1\tau_1\bigl(\hat{\boldsymbol e}_u, \alpha X(V_h)\bigr)_K,
\]
and since $\varphi_1\tau_1 \le \varphi_1^{1/2}\tau_1^{1/2}$ (from
$\varphi_1\tau_1 \le 1$) with $\varphi_1^{1/2} =
\sqrt{c_1}\,(\alpha_K/h_K)\,\tau_{2,K}^{1/2}$,
\[
  \Bigl| \sum_K \varphi_1\tau_1\bigl(\hat{\boldsymbol e}_u,
  \alpha X(V_h)\bigr)_K \Bigr|
  \le \sqrt{c_1}\sum_K \frac{\alpha_K}{h_K}\,\tau_{2,K}^{1/2}\,\Eint{K}{\bu}\;
  \tau_{1,K}^{1/2}\normK{\alpha X(V_h)}
  \le C\,\EF(h)\,\tnorm{V_h}_{\mathrm{S}},
\]
again Damk\"ohler-free. Both dangerous first-order Galerkin pairings are thus
neutralized by the reactive parts of the ASGS trial and test slots; this is where
the full-residual structure pays, and precisely what the OSGS annihilations forgo
(cf.\ \cref{rem:inert}).

\smallskip\noindent\emph{Stabilization part: inventory.}
Expanding \cref{eq:BSdef} at $(\hat E, V_h)$, the momentum stabilization is the sum
over $K$ of $\tau_{1,K}$ times the products of the test factors
$\{\alpha X(V_h),\ 2\nabla\cdot(\alpha\nu\ViscProj\nabla\bv_h),\ -\sigma\bv_h\}$
with the trial factors
$\{\alpha X(\hat E),\ -2\nabla\cdot(\alpha\nu\ViscProj\nabla\hat{\boldsymbol e}_u),\
+\sigma\hat{\boldsymbol e}_u\}$; the reactive--reactive product and the
$\bigl(\alpha X(V_h), \sigma\hat{\boldsymbol e}_u\bigr)$ product were consumed by
the two telescopings above. The mass stabilization contributes the products of
$\{\nabla\cdot(\alpha\bv_h),\ -\varepsilon q_h\}$ with
$\{\varepsilon\hat e_p,\ \nabla\cdot(\alpha\hat{\boldsymbol e}_u)\}$. We bound the
factors elementwise once and multiply.

\smallskip\noindent\emph{Trial factors.}
\begin{align}
  \tau_{1,K}^{1/2}\normK{\alpha X(\hat E)}
  &\le C\,\frac{\alpha_K}{h_K}\bigl( \tau_{2,K}^{1/2}\Eint{K}{\bu}
  + \tau_{1,K}^{1/2}\Eint{K}{p} \bigr)
  && \text{(as in (I6) of~\cite{OSGSnote})},
  \label{eq:TrialX}\\
  \tau_{1,K}^{1/2}\normK{2\nabla\cdot(\alpha\nu\ViscProj\nabla\hat{\boldsymbol e}_u)}
  &\le C(1+C_\nabla)\,\frac{\alpha_K}{h_K}\,\tau_{2,K}^{1/2}\Eint{K}{\bu}
  && \text{(Damk\"ohler-free)},
  \label{eq:TrialVisc}\\
  \tau_{1,K}^{1/2}\normK{\sigma\hat{\boldsymbol e}_u}
  \le \sigma^{1/2}\normK{\hat{\boldsymbol e}_u}
  &\le C\,\Dah^{1/2}\,\frac{\alpha_K}{h_K}\,\tau_{2,K}^{1/2}\Eint{K}{\bu}
  && \text{(penalized; (P6))},
  \label{eq:TrialReac}\\
  \tau_{2,K}^{1/2}\normK{\nabla\cdot(\alpha\hat{\boldsymbol e}_u)}
  &\le C(1+C_\nabla)\,\frac{\alpha_K}{h_K}\,\tau_{2,K}^{1/2}\Eint{K}{\bu}
  && \text{(as in (I7) of~\cite{OSGSnote})},
  \label{eq:TrialDiv}\\
  \tau_{2,K}^{1/2}\normK{\varepsilon \hat e_p}
  \le C_2\,\tau_{2,K}^{-1/2}\normK{\hat e_p}
  &\le C\,(c_1+\Dah)^{1/2}\,\frac{\alpha_K}{h_K}\,\tau_{1,K}^{1/2}\Eint{K}{p}
  && \text{(penalized; (P1), (P5))}.
  \label{eq:TrialEps}
\end{align}
For \cref{eq:TrialVisc}: expanding the divergence,
$\normK{\nabla\cdot(\alpha\nu\ViscProj\nabla\hat{\boldsymbol e}_u)}
\le \nu\bigl( \alpha_K \normK{\nabla(\ViscProj\nabla\hat{\boldsymbol e}_u)}
+ \norm{\nabla\alpha}_{L^\infty(K)}\normK{\ViscProj\nabla\hat{\boldsymbol e}_u}
\bigr)
\le C(1+C_\nabla)\,\nu\alpha_K\, h_K^{k_u-1}\norm{\bu}_{H^{k_u+1}(K)}$
(the interpolation estimate with $m = 2$ and the resolution condition; $\ViscProj$
has constant coefficients), and then
$\tau_{1,K}^{1/2}\nu\alpha_K h_K^{k_u-1}
\le \alpha_K^{1/2}\nu^{1/2}h_K^{k_u}/\sqrt{c_1}
\le \alpha_K\tau_{2,K}^{1/2}h_K^{k_u}/\sqrt{c_1}$
via $\tau_1 \le h^2/(c_1\alpha_K\nu)$ and $\nu \le \alpha_K\tau_{2,K}$. Note that
the trial side never needs an inverse estimate: $\hat E$ is an interpolation error,
not a finite element function.

\smallskip\noindent\emph{Test factors.}
\begin{align}
  \tau_{1,K}^{1/2}\normK{\alpha X(V_h)},\quad
  \tau_{2,K}^{1/2}\normK{\nabla\cdot(\alpha\bv_h)}
  &\;\le\; \tnorm{V_h}_{\mathrm{S}}\text{-summands},
  \label{eq:TestFirst}\\
  \tau_{1,K}^{1/2}\normK{2\nabla\cdot(\alpha\nu\ViscProj\nabla\bv_h)}
  &\le \frac{2\,\delta_\alpha^{1/2} M_2}{\gamma}\,
  \nu^{1/2}\bigl\lVert \alpha^{1/2}\ViscProj\nabla\bv_h \bigr\rVert_K
  && \text{(Damk\"ohler-free)},
  \label{eq:TestVisc}\\
  \tau_{1,K}^{1/2}\normK{\sigma\bv_h}
  = (\sigma^2\tau_{1,K})^{1/2}\normK{\bv_h}
  &\le \Bigl(\frac{\Dah}{c_1}\Bigr)^{1/2} \st^{1/2}\normK{\bv_h}
  && \text{(penalized; \cref{eq:ThreeIdentities})},
  \label{eq:TestReac}\\
  \tau_{2,K}^{1/2}\normK{\varepsilon q_h}
  \le (\varepsilon\tau_{2,K})^{1/2}\,\varepsilon^{1/2}\normK{q_h}
  &\le C_2^{1/2}\,\varepsilon^{1/2}\normK{q_h}
  && \text{(Damk\"ohler-free)}.
  \label{eq:TestEps}
\end{align}
For \cref{eq:TestVisc}, the same divergence expansion and the elementwise bound
$\normK{\ViscProj\nabla\bv_h} \le \alpha_{0,K}^{-1/2}
\normK{\alpha^{1/2}\ViscProj\nabla\bv_h}$ give the companion's enlarged inverse
estimate with constant $\delta_\alpha^{1/2}M_2 C_{\mathrm{inv}}$, after which (K1)
applies; this is the one place the enlarged constant is needed, and it is needed
because the ASGS \emph{test slot} genuinely contains the viscous operator---an
irreducible difference from the OSGS analysis.

\smallskip\noindent\emph{Assembly.}
Multiplying test and trial factors, summing over $K$ by Cauchy--Schwarz, exactly
three families carry a Damk\"ohler factor: (F2)
$\bigl(2\nabla\cdot(\alpha\nu\ViscProj\nabla\bv_h),
\sigma\hat{\boldsymbol e}_u\bigr)_{\tau_1}$, penalized through the trial factor
\cref{eq:TrialReac}; and (F3) $\bigl(-\sigma\bv_h, \alpha X(\hat E)\bigr)_{\tau_1}$
and (F4) $\bigl(-\sigma\bv_h,
-2\nabla\cdot(\alpha\nu\ViscProj\nabla\hat{\boldsymbol e}_u)\bigr)_{\tau_1}$,
penalized through the test factor \cref{eq:TestReac} (the former (F1) family,
$(\alpha X(V_h), \sigma\hat{\boldsymbol e}_u)_{\tau_1}$, having been telescoped
away). The
$\varepsilon$-products carry at most the $(c_1+\Dah)^{1/2}$ of \cref{eq:TrialEps}.
All remaining products are Damk\"ohler-free. Every product is bounded by
$C(1+\Dah)^{1/2}$ times the corresponding term of $\EF(h)$ times a summand of
$\tnorm{V_h}_{\mathrm{S}}$, and \cref{eq:ASGSInterpBound} follows.
\end{proof}

\begin{theorem}[Unified penalized convergence, ASGS instantiation]
\label{th:unifiedASGS}
Let $U$ solve the linearized problem and $U_h \in \Xhz$ the ASGS discrete problem
of~\cite{companion}. Under the hypotheses stated at the head of this section
(in particular: no jump condition, no hypothesis (A7)),
\begin{equation}
  \label{eq:UnifiedASGS}
  \tnorm{U - U_h}_{\mathrm{S}}
  \;\le\;
  C \sum_K \frac{\alpha_K}{h_K}\,\bigl(1+\Dah\bigr)^{1/2}
  \Bigl( \tau_{2,K}^{1/2}\Eint{K}{\bu} + \tau_{1,K}^{1/2}\Eint{K}{p} \Bigr).
\end{equation}
\end{theorem}

\begin{proof}
The skeleton of~\cite[Theorem~6.4]{OSGSnote}, with two substitutions.
\emph{Stability}: the coercivity lemma of~\cite{companion} gives
$\Bs(\ba; W_h, W_h) \ge \beta\,\tnorm{W_h}_{\mathrm{S}}^2$ for $W_h \in \Xhz$; this
is the degenerate (empty test enrichment) instance of the inf--sup template, cf.\
\cref{rem:deficits}. \emph{Consistency}: the ASGS method is variationally
consistent---for the exact solution, $\alpha X(U) - G(\bu) = \bff$ and
$\varepsilon p + \nabla\cdot(\alpha\bu) = 0$, so the stabilization terms
of~\cref{eq:BSdef} evaluate against the data exactly as in the companion's
right-hand side, whence $\Bs(\ba; U - U_h, V_h) = 0$ for all $V_h \in \Xhz$; the
best-approximation consistency lemma of~\cite{OSGSnote} is replaced by this
identity. Then, with $\eta_h = \hat U_h - U_h$,
\[
  \beta\,\tnorm{\eta_h}_{\mathrm{S}}
  \le \frac{\Bs(\ba;\eta_h,\eta_h)}{\tnorm{\eta_h}_{\mathrm{S}}}
  = \frac{-\Bs(\ba;\hat E,\eta_h)}{\tnorm{\eta_h}_{\mathrm{S}}}
  \le C\,\EF(h)
\]
by \cref{lem:asgsinterp}, and
$\tnorm{\hat E}_{\mathrm{S}} \le \tnorm{\hat E} \le C\,\EF(h)$
by~\cite[Lemma~6.5]{OSGSnote} and the termwise norm comparison. The triangle
inequality concludes.
\end{proof}

\begin{remark}[Where the penalty actually resides; further repairs]
\label{rem:sharper}
The three surviving penalized families are not equally stubborn. Family (F2) can be
repaired outright, with no new hypothesis: elementwise integration by parts gives
\begin{align*}
  \sigma\tau_{1,K}\bigl(2\nabla\cdot(\alpha\nu\ViscProj\nabla\bv_h),
  \hat{\boldsymbol e}_u\bigr)_K
  ={}& -2\sigma\tau_{1,K}\bigl(\alpha\nu\ViscProj\nabla\bv_h,
  \ViscProj\nabla\hat{\boldsymbol e}_u\bigr)_K \\
  &+ \sigma\tau_{1,K}\int_{\partial K} 2\alpha\nu\,
  (\ViscProj\nabla\bv_h\,\bn)\cdot\hat{\boldsymbol e}_u\,\mathrm{d}\Gamma;
\end{align*}
using $\sigma\tau_1 \le 1$ and $\nu \le \alpha_K\tau_{2,K}$, the volume term is
bounded by
$C\delta_\alpha^{1/2}\,(\alpha_K/h_K)\tau_{2,K}^{1/2}\Eint{K}{\bu}\cdot
\nu^{1/2}\lVert\alpha^{1/2}\ViscProj\nabla\bv_h\rVert_K$, and the face terms,
estimated per element by a discrete trace inequality on $\ViscProj\nabla\bv_h$ and a
continuous one on $\hat{\boldsymbol e}_u$, carry the same weight; both are
Damk\"ohler-free, and no jump quantity arises because each face integral is bounded
directly rather than assembled facewise. The genuinely delicate families are (F3),
(F4) and the pressure/$\varepsilon$ conversions of \cref{eq:TrialEps} and (I5):
these are exactly the terms on which the companion's proof spends its facewise
$[\tau_1]$-estimates (its terms [I]+[IV] and [III], in the notation of the
companion's appendix), and they can be made Damk\"ohler-free by two routes: (a) the
companion's---elementwise integration by parts, facewise assembly into
$[\tau_1]$-jumps, and the companion's jump lemma; or (b) a regime-splitting
route---partition the elements by $\sigma\tau_1 \gtrless 1/2$, bound the
reaction-dominated side directly (there $\st \ge \varphi_1/2$, so
$\st^{1/2}$-pairings saturate the $\varphi_1$-scale), integrate by parts only on the
complement (where $\Dah \lesssim c_1$, so every conversion is bounded), and estimate
the interface faces by trace inequalities. We have verified the weights of route (b)
family by family; both routes rest on the same facewise comparability of $\varphi_1$,
which \cref{lem:jumpauto} below shows is implied by the standing hypotheses. We have
not written route (b) out in full and do not claim it as a theorem; its interest is
conceptual---it indicates that the jump \emph{machinery}, not merely the jump
\emph{hypothesis}, is replaceable.
\end{remark}

\section{Comparison, corollaries and the residual asymmetries}
\label{sec:comparison}

\begin{remark}[What is shared and what is not]
\label{rem:shared}
The unified pair---\cref{th:unifiedASGS} here and
\cite[Theorem~6.4]{OSGSnote}---now has the following anatomy:
\begin{center}
\small
\begin{tabular}{lll}
\toprule
component & ASGS instantiation & OSGS instantiation \\
\midrule
hypotheses & (A1)--(A6), (A8) & (A1)--(A8) \\
extra hypothesis & enlarged inverse const.\ (automatic) & (A7); not automatic \\
parameter toolbox & \multicolumn{2}{c}{(P1)--(P6), (K1)--(K3): shared} \\
stability & coercivity (empty enrichment) & inf--sup (projection enrichment) \\
consistency & exact ($\equiv 0$) & best approximation
(\cite[Lemma~6.1]{OSGSnote}) \\
interpolation core & (I1), (I2), (I4), (I5) + telescoped (I3$'$)
& (I1)--(I5) \\
stabilization terms & inventory of \cref{lem:asgsinterp} & nonexpansiveness
(2 lines) \\
norm & $\tnorm{\cdot}_{\mathrm{S}}$ ($\st$; forced by \cref{prop:negative}) &
$\tnorm{\cdot}$ (full $\sigma$) \\
result & $\EF(h)$, i.e., $(1+\Dah)^{1/2}\EF_0(h)$ &
$\EF(h)$, identical \\
$\varphi_1$-face comparability & $[\tau_1]$-face estimates
(\cref{lem:jumpauto}) & patch equiv., smoothing \\
\bottomrule
\end{tabular}
\end{center}
A single parametric write-up covering both columns is entirely feasible; the
method-specific residue is the two extra-hypothesis cells, the consistency cell,
and the stabilization-term cell.
\end{remark}

\begin{remark}[$L^2$ velocity: both methods, both routes]
\label{rem:Ltwo}
The companion's sharp Theorem~5.3 yields, in the reaction-dominated regime, using
$\st \eqsim c_1\alpha_K^2\tau_{2,K}/h_K^2$ there,
$\norm{\bu-\bu_h}_{\mathrm{ASGS}} \le \st^{-1/2}\,C\,\EF_0(h) \le
C\bigl( \Eint{K}{\bu}\text{-terms} + (\tau_1/\tau_2)^{1/2}\Eint{K}{p}\text{-terms}
\bigr)$: optimal order with a $\sigma$-uniform constant. The OSGS corollary
of~\cite[Corollary~6.6]{OSGSnote} gives
$\norm{\bu-\bu_h}_{\mathrm{OSGS}} \le C\sum_K(1+\Dah^{-1})^{1/2}
[(1+\tfrac{c_2}{c_1}\Reh)^{1/2}\Eint{K}{\bu} +
(\alpha_K/(\sigma\nu))^{1/2}\Eint{K}{p}]$: also optimal order,
$\sigma$-robust. So both analyses predict indistinguishable $L^2$ velocity behavior
at strong reaction, which is exactly what the tables of~\cite{companion} show; the
Damk\"ohler asymmetry lives entirely in the $H^1$-type components, again as
observed.
\end{remark}

\begin{remark}[A third instantiation, and why the full-residual ASGS resists the
template]
\label{rem:thirdinstance}
The Method-I template of~\cite{Codina2008} with the projection set to the identity
complement removed---i.e., stabilizing the plain first-order quantities
$(\alpha X(U_h), \alpha X(V_h))_{\tau_1}$ and
$(\nabla\cdot(\alpha\bu_h), \nabla\cdot(\alpha\bv_h))_{\tau_2}$ without any
projection---is the porous version of the term-by-term method
of~\cite{Chacon1998}. It fits the unified skeleton perfectly (trivial stability
enrichment, no (A7)), but its consistency error is
$(\alpha X(U), \alpha X(V_h))_{\tau_1}$-sized, i.e., $O(\tau_1^{1/2})$ times an
$O(1)$ residual, which restricts it to first-order elements, as~\cite{Codina2008}
already notes. This clarifies the obstruction from the other side: the
\emph{full-residual} structure of the companion's ASGS is exactly what buys exact
consistency for all orders---and exactly what installs the viscous and reactive
operators in both slots of the quadratic form, forcing the renormalization of
\cref{prop:negative} and the enlarged inverse constant. One cannot keep the asset
and drop the liabilities.
\end{remark}

\begin{question}[updated]
\label{q:open}
The original question asked whether the family
$(\alpha X(V_h), \sigma\hat{\boldsymbol e}_u)_{\tau_1}$ admits a bound that is
simultaneously Damk\"ohler-uniform and free of conditions on the jumps of
$\tau_1$.
It is now largely resolved, by dissolution rather than by answer: (i) that family
telescopes against the Galerkin first-order pairing at zero cost
(\cref{lem:asgsinterp}, second telescoping), with no integration by parts and no
face terms; (ii) the families that do route through facewise machinery ((F3), (F4)
and the pressure/$\varepsilon$ conversions) admit, besides the companion's
$[\tau_1]$-route, the regime-splitting route sketched in \cref{rem:sharper}; and
(iii) both routes rest only on the facewise comparability of $\varphi_1$, which
\cref{lem:jumpauto} shows is implied by the standing hypotheses. The jump condition
was therefore never an independent assumption, as the companion's own remark on it
anticipated. What remains open is bookkeeping, not principle: a full write-up of the
regime-splitting continuity lemma (which would make the companion's appendix
independent of the explicit jump estimates), and, on the impossibility side, whether
the Damk\"ohler factor can be removed from the OSGS bound in norms \emph{weaker}
than the strong one---\cref{prop:floor} settles the strong norm, and the
missing-telescoping-partner obstruction blocks all continuity-type arguments in the
$\st$-norm, but duality-type arguments are not excluded.
\end{question}

\section{The jump condition examined: automatic for ASGS, inert for OSGS}
\label{sec:jump}

The companion's continuity lemma assumes, when $\sigma > 0$, the facewise
comparability
\begin{equation}\label{eq:Hjump}
  c_J\,\varphi_{1,K_1} \;\le\; \varphi_{1,K_2} \;\le\; c_J'\,\varphi_{1,K_1},
  \qquad \varphi_{1,K} \coloneqq \alpha_K\,\tns^{-1},
  \qquad 0 < c_J \le 1 \le c_J',
\end{equation}
for every pair of elements sharing a face (its hypothesis H\textsubscript{jump},
the porous analogue of condition~(39) in~\cite{Codina2001}). A natural question,
prompted by the comparison of the two analyses: is this an extra assumption whose
availability would \emph{also} improve the OSGS estimate---so that listing its
absence as an OSGS feature misreads the ledger? The answer is that the condition is
(i) implied by the standing hypotheses, (ii) identical in content to the patch
equivalence on which the OSGS analysis already rests, and (iii) without any point of
application in the OSGS analysis; and that the OSGS penalty one might have hoped it
removes is unremovable by \emph{any} hypothesis. We take these in turn.

\begin{lemma}[The jump condition is implied by the standing hypotheses]
\label{lem:jumpauto}
Assume the nondegeneracy hypothesis (A1), the porosity bounds and resolution
condition (A3) with elementwise contrast $\delta_\alpha$, and continuity of the
advection field, $\ba \in C^0(\overline\Omega)^d$, with modulus of continuity
$\omega_{\ba}$; let $\nu > 0$. Then there is $h_0 > 0$, depending only on $\nu$ and
$\omega_{\ba}$, such that \cref{eq:Hjump} holds for all $h \le h_0$ with constants
$c_J, c_J'$ depending only on the mesh-regularity constants, $\delta_\alpha$, $c_1$
and $c_2$. Since $\sigma$ cancels in the jump of $\tau_1^{-1} = \varphi_1 + \sigma$,
the condition constrains only the data and the mesh, never the reaction.
\end{lemma}

\begin{proof}
Let $K_1, K_2$ share the face $F$. \emph{Porosity factors:} since
$\alpha_{0,K_1} \le \inf_F \alpha$ and $\inf_F\alpha \le \alpha_{K_2}$,
\[
  \alpha_{K_1} \le \delta_\alpha\,\alpha_{0,K_1}
  \le \delta_\alpha \inf_F \alpha \le \delta_\alpha\,\alpha_{K_2},
\]
and symmetrically; this is the argument of the companion's own remark on the
condition. \emph{Mesh factors:} $h_{K_1} \eqsim h_{K_2}$ with constants from (A1).
\emph{Advection factors:} $\bigl|\,|\ba|_{\infty,K_1} - |\ba|_{\infty,K_2}\,\bigr|
\le \omega_{\ba}(C h)$ since $K_1 \cup K_2$ has diameter $\le Ch$, whence
\[
  \frac{c_1\nu}{h_{K_1}^2} + \frac{c_2|\ba|_{\infty,K_1}}{h_{K_1}}
  \;\le\; C\Bigl( \frac{c_1\nu}{h_{K_2}^2}
  + \frac{c_2|\ba|_{\infty,K_2}}{h_{K_2}} \Bigr)
  + \frac{C\,c_2\,\omega_{\ba}(Ch)}{h_{K_2}},
\]
and the last term is absorbed by the viscous floor:
$c_2\omega_{\ba}(Ch)/h_{K_2} = \bigl[c_2\,\omega_{\ba}(Ch)\,h_{K_2}/(c_1\nu)\bigr]
\cdot c_1\nu/h_{K_2}^2$, whose bracket is $\le 1$ for
$h \le h_0(\nu, \omega_{\ba})$. Multiplying the three comparabilities gives
\cref{eq:Hjump}. Two observations. First, the $h_0(\nu)$-dependence is exactly that
of the smoothing function $\psi(h)$ in the OSGS stability proof, which controls the
oscillation of the same quantity, by the same viscous floor, for the same reason.
Second, \cref{eq:Hjump} restricted to face-neighbours is precisely the comparability
content of the patch-equivalence lemma~\cite[Lemma~4.2]{OSGSnote} that the OSGS
interpolation and smoothing estimates consume. The condition is thus not an
independent assumption but a restatement, at faces, of regularity the standing
hypotheses already provide.
\end{proof}

\begin{remark}[No point of application in the OSGS analysis]
\label{rem:inert}
Suppose one grants \cref{eq:Hjump} to the OSGS analysis. Nothing improves, for a
structural reason. In the companion's continuity proof the condition enters through
a single consequence: the facewise estimate
$\sigma\,\bigl|[\tau_1]\bigr| \le C\,
(\st\tau_1|_{K_1})^{1/2}(\st\tau_1|_{K_2})^{1/2}$, which prices the interior-face
terms generated by \emph{elementwise} integration by parts of the $\sigma$-weighted
stabilization pairings---the fee for executing the telescopings
$1 - \sigma\tau_1 = \varphi_1\tau_1$ on its terms [I]+[IV] and [III] (in the
notation of the companion's appendix; the third telescoping, [II]+[V], is jump-free
algebra, and is the one imported into \cref{lem:asgsinterp}). The pairings being
telescoped exist because the ASGS test slot contains $-\sigma\bv_h$ and its residual
slot $+\sigma\bu_h$. The OSGS projection annihilates precisely these
(the reactive parts of its residual and test slots have no fluctuation, cf.\ the
form \cref{eq:Bodef}): there is no partner, hence no telescoping, hence no
elementwise integration by parts, hence no $[\tau_1]$ face term anywhere in the OSGS
error analysis. The Damk\"ohler losses there sit instead in \emph{Galerkin}
pairings---$\sigma(\hat{\boldsymbol e}_u, \bv_h)$, the $\sigma^{1/2}$-part of the
$\alpha X$-pairing, the $\tau_2^{-1/2}$ pressure weight---which carry no
$\tau_1$-weight and, in the reactive case, no derivative structure to move. A
hypothesis about $[\tau_1]$ has nothing to act on. Note finally that the OSGS proof
already \emph{uses} the comparability content of \cref{eq:Hjump}, as patch
equivalence, in its quasi-interpolation and smoothing lemmas; what it lacks is not
the hypothesis but the telescoping partners.
\end{remark}

\begin{proposition}[The strong-norm penalty is forced by the norm]
\label{prop:floor}
Let $\ba \equiv \boldsymbol 0$, $\alpha \equiv 1$, $\varepsilon = 0$, on
quasi-uniform meshes, and fix
$\bu \in C^\infty(\overline\Omega)^d \cap H^1_0(\Omega)^d$ with
$\nabla\cdot\bu = 0$ and $D^{k_u+1}\bu \not\equiv 0$; set $p \coloneqq 0$ and
$\bff \coloneqq \sigma\bu - 2\nu\,\nabla\cdot(\ViscProj\nabla\bu)$, so that
$U = [\bu; 0]$ solves the linearized problem of~\cite{OSGSnote} for every
$\sigma, \nu > 0$. Then there is $c_{\bu} > 0$, independent of $h$, $\sigma$ and
$\nu$, such that for \emph{every} $V_h \in \Xhz$---in particular for the discrete
solution of any method whatsoever---
\[
  \tnorm{U - V_h} \;\ge\; \norm{\sigma^{1/2}(\bu - \bv_h)}
  \;\ge\; c_{\bu}\,\sigma^{1/2} h^{k_u+1}
  \;=\; \frac{c_{\bu}}{\norm{\bu}_{H^{k_u+1}(\Omega)}}\;
  \mathrm{Da}_h^{1/2}\;\EF_0^{(2)}(h),
\]
where $\EF_0^{(2)}(h) \coloneqq \bigl( \sum_K \bigl[(\alpha_K/h_K)\,
\tau_{2,K}^{1/2}\,\Eint{K}{\bu}\bigr]^2 \bigr)^{1/2} = \nu^{1/2} h^{k_u}
\norm{\bu}_{H^{k_u+1}(\Omega)}$ in this setting. Consequently no a priori estimate
$\tnorm{U - U_h} \le C\,\EF_0^{(2)}(h)$ with $C$ independent of $\mathrm{Da}_h$ can
hold for any discretization; and since $\sigma$ is a free parameter, taking $\sigma$
large at fixed $h$ extends the conclusion to the summed error function $\EF_0(h)
\ge \EF_0^{(2)}(h)$. The factor $(1+\Dah)^{1/2}$ in the OSGS theorem
of~\cite{OSGSnote} is therefore optimal up to constants: it measures the norm, not
the proof. This complements \cref{prop:negative}: that result showed the ASGS form
cannot \emph{reach} the strong norm; this one shows nothing can \emph{beat} the
strong norm's own best-approximation scaling.
\end{proposition}

\begin{proof}
Only the best-approximation lower bound requires proof. Discrete velocities are
continuous piecewise polynomials of degree $k_u$, so
$\operatorname{dist}_{L^2}(\bu, V_{h0}) \ge
\operatorname{dist}_{L^2}\bigl(\bu, \mathbb P_{k_u}(\mathcal T_h)\bigr)$, the
unconstrained piecewise-polynomial distance. Pick a ball $\omega \subset \Omega$
and $\delta > 0$ with $|D^{k_u+1}\bu| \ge \delta$ on $\omega$ (possible:
$D^{k_u+1}\bu$ is continuous and not identically zero). For $K \subset \omega$, let
$T_K$ be the Taylor polynomial of $\bu$ of degree $k_u{+}1$ at the barycenter and
$H_K$ its homogeneous part of degree $k_u{+}1$. Then
\[
  \operatorname{dist}_{L^2(K)}(\bu, \mathbb P_{k_u})
  \;\ge\; \bigl\lVert (I - \Pi_{\mathbb P_{k_u}(K)})\,H_K \bigr\rVert_{L^2(K)}
  - \norm{\bu - T_K}_{L^2(K)}
  \;\ge\; c\,\delta\, h_K^{k_u+1}\,|K|^{1/2} - C_{\bu}\, h_K^{k_u+2}\,|K|^{1/2},
\]
where the first bound holds because
$H \mapsto \lVert (I - \Pi_{\mathbb P_{k_u}}) H \rVert_{L^2(\widehat K)}$ is a norm
on the homogeneous polynomials of degree $k_u{+}1$ (which meet
$\mathbb P_{k_u}$ only at $0$), hence equivalent to the coefficient norm on that
finite-dimensional space, uniformly over shape-regular elements after affine
scaling; and the second is the Taylor remainder. Squaring, summing over
$K \subset \omega$, and taking $h$ small enough that the remainder halves the
leading term and the elements contained in $\omega$ cover half its measure,
\[
  \operatorname{dist}_{L^2}\bigl(\bu, \mathbb P_{k_u}(\mathcal T_h)\bigr)^2
  \;\ge\; \tfrac14\, c^2\delta^2\, |\omega|\, h^{2(k_u+1)}
  \;\eqqcolon\; c_{\bu}^2\, h^{2(k_u+1)}.
\]
The identifications $\tau_{2,K} = \nu$, $\alpha_K = 1$ (so that
$(\alpha_K/h_K)\tau_{2,K}^{1/2}\Eint{K}{\bu} = \nu^{1/2}h^{k_u}
\norm{\bu}_{H^{k_u+1}(K)}$ up to the quasi-uniformity constants) and
$\sigma^{1/2}h^{k_u+1} = \mathrm{Da}_h^{1/2}\,\nu^{1/2}h^{k_u}$ complete the
computation.
\end{proof}

\begin{remark}[The corrected hypothesis ledger]
\label{rem:ledger}
\Cref{lem:jumpauto,rem:inert,prop:floor} correct a framing that the comparison
table of \cref{rem:shared}, and the notes~\cite{OSGSnote}, might otherwise invite.
``No jump condition needed'' is not an advantage of the OSGS analysis: the condition
costs nothing (it is implied by the standing hypotheses, exactly as the patch
equivalence and smoothing estimates that the OSGS proof itself consumes), and having
it would buy nothing (it has no point of application, and the dominant penalty is
norm-forced). Symmetrically, the enlarged inverse-estimate constant of the ASGS
analysis is $\delta_\alpha^{1/2} M_2\, C_{\mathrm{inv}}$, derived from the same
resolution condition. Both analyses therefore rest on the same automatic regularity
substrate; the ASGS proof converts it into Damk\"ohler-free sharpness through the
three telescopings, and the OSGS proof cannot. The single assumption in either
analysis that is \emph{not} implied by the standing data and mesh hypotheses is the
projection-stability condition (A7), on the OSGS side, whose verification is a
genuine property of the finite element spaces
(\cite{CodinaBlasco1997}-type). If the hypothesis ledger tilts at all, it tilts
toward the ASGS analysis.
\end{remark}

\section{Recommendation for the manuscript}
\label{sec:recommendation}

Three options, in increasing order of disruption. \emph{Option 1 (recommended for
the present submission):} keep both existing proofs exactly as they are, and add one
short remark after the ASGS convergence theorem noting (i) that the OSGS analysis
of the companion note follows the same architecture with the projection-adapted
test function, (ii) the dichotomy of \cref{prop:negative,rem:symmetry}
($\st$ sharp for ASGS, full $\sigma$ with a $(1+\Dah)^{1/2}$ constant for OSGS),
(iii) that this dichotomy is consistent with the reaction-dominated
observations of the numerical section, and (iv) that the jump condition is implied
by the standing hypotheses (\cref{lem:jumpauto}), so the ASGS theorem's effective
hypothesis set is no larger than the standing one. Cost: a paragraph. \emph{Option 2:}
restructure the analysis section as a single parametric proof
(\cref{rem:shared}'s table as the organizing device) with the sharp
Damk\"ohler-free ASGS bound retained as a method-specific refinement via the jump
condition. This is the intellectually cleanest presentation but is a
major-revision-scale rewrite and moves the jump condition to a more prominent
position than it currently occupies. \emph{Option 3:} unified analysis in
supplementary material, main text unchanged. Our assessment: the unified theorem's
value is explanatory (it certifies the mechanism behind the observed OSGS/ASGS
differences and removes the jump condition at a quantified price), not an
improvement of any individual estimate; that argues for Option~1 now, with
Options~2--3 as candidates for a follow-up methods paper if the open question above
is resolved.

\begin{thebibliography}{9}

\bibitem{companion}
G.~Casas, J.~Gonz\'alez-Us\'ua, R.~Codina, I.~de-Pouplana,
\emph{A stabilized finite element method for incompressible, inertial flows in
inhomogeneous porous media}, companion manuscript (submitted).

\bibitem{OSGSnote}
\emph{Stability and a priori error analysis of the OSGS variant of a stabilized
finite element method for incompressible, inertial flows in inhomogeneous porous
media}, companion note to~\cite{companion} (working document,
\texttt{osgs\_convergence.pdf}).

\bibitem{Codina2008}
R.~Codina,
\emph{Analysis of a stabilized finite element approximation of the Oseen equations
using orthogonal subscales},
Applied Numerical Mathematics 58 (2008) 264--283.

\bibitem{Codina2001}
R.~Codina,
\emph{A stabilized finite element method for generalized stationary incompressible
flows},
Computer Methods in Applied Mechanics and Engineering 190 (2001) 2681--2706.

\bibitem{CodinaBlasco1997}
R.~Codina, J.~Blasco,
\emph{A finite element formulation for the Stokes problem allowing equal
velocity--pressure interpolation},
Computer Methods in Applied Mechanics and Engineering 143 (1997) 373--391.

\bibitem{Chacon1998}
T.~Chac\'on Rebollo,
\emph{A term by term stabilization algorithm for finite element solution of
incompressible flow problems},
Numerische Mathematik 79 (1998) 283--319.

\bibitem{BrennerScott}
S.~C. Brenner, L.~R. Scott,
\emph{The Mathematical Theory of Finite Element Methods},
Springer, New York, 2008.

\end{thebibliography}

\end{document}
